\documentclass[../../main.tex]{subfiles}
\graphicspath{{../../image/}{../../figure/}}
\begin{document}

\section{主要结果}\label{section:ch15-主要结果}

我们考虑 Cauchy 问题
\begin{align}
\partial_t U(x,t)+\partial_x F(U(x,t))=0,\quad&-\infty<x<\infty,\ 0<t<\infty,
\label{align:ch15-1511}\\
U(x,0)=U_0(x),\quad&-\infty<x<\infty,
\label{align:ch15-1512}
\end{align}
其中守恒律系统在 $\mathbb{R}^n$ 中以某状态 $U^*$ 为中心的球 $\mathscr{O}$ 内严格双曲，初始数据 $U_0$ 是 $(-\infty,\infty)$ 上的有界变差函数，且满足 $U_0(-\infty)=U^*$。

目标是把\eqref{align:ch15-1511}、\eqref{align:ch15-1512}的 $BV$ 解构造为抛物系统
\begin{align}
\partial_t U(x,t)+\partial_x F(U(x,t))=\mu\,\partial_x^2U(x,t),\quad-\infty<x<\infty,\ 0<t<\infty,
\label{align:ch15-1513}
\end{align}
在同一初始条件\eqref{align:ch15-1512}下的解当 $\mu\downarrow0$ 时的极限。这将通过下述结果完成。

\begin{theorem}[消失黏性主定理]\label{theorem:ch15-消失黏性主定理}
存在 $\delta>0$，使得当
\begin{align}
TV_{(-\infty,\infty)}U_0(\cdot)<\delta
\label{align:ch15-1514}
\end{align}
时，对某个正常数 $a$ 与 $b$，下述结论成立：
\begin{enumerate}
\item\label{item:ch15-151-a} 对任意 $\mu>0$，存在\eqref{align:ch15-1513}、\eqref{align:ch15-1512}的一个古典解 $U_\mu$，满足
\begin{align}
TV_{(-\infty,\infty)}U_\mu(\cdot,t)&\leqslant a\,TV_{(-\infty,\infty)}U_0(\cdot),\qquad t\in(0,\infty),
\label{align:ch15-1515}\\
\|U_\mu(\cdot,t)-U_\mu(\cdot,\tau)\|_{L^1(-\infty,\infty)}&\leqslant b\bigl(|t-\tau|+|\sqrt{\mu t}-\sqrt{\mu\tau}|\bigr),\qquad \tau,t\in(0,\infty).
\label{align:ch15-1516}
\end{align}
\item\label{item:ch15-151-b} 若 $\bar U_\mu$ 表示\eqref{align:ch15-1513}以 $\bar U_0$ 为初值的解，其中 $TV_{(-\infty,\infty)}\bar U_0(\cdot)<\delta$，$\bar U_0(-\infty)=U^*$，且 $U_0-\bar U_0\in L^1(-\infty,\infty)$，则
\begin{align}
\|U_\mu(\cdot,t)-\bar U_\mu(\cdot,t)\|_{L^1(-\infty,\infty)}\leqslant a\|U_0(\cdot)-\bar U_0(\cdot)\|_{L^1(-\infty,\infty)},\quad t\in(0,\infty).
\label{align:ch15-1517}
\end{align}
\item\label{item:ch15-151-c} 当 $\mu\downarrow0$ 时，$\{U_\mu\}$ 在 $L^1_{\mathrm{loc}}$ 中收敛到\eqref{align:ch15-1511}、\eqref{align:ch15-1512}的一个 $BV$ 解 $U$，它继承稳定性估计\eqref{align:ch15-1515}、\eqref{align:ch15-1516}与\eqref{align:ch15-1517}，即
\begin{align}
TV_{(-\infty,\infty)}U(\cdot,t)&\leqslant a\,TV_{(-\infty,\infty)}U_0(\cdot),\qquad t\in(0,\infty),
\label{align:ch15-1518}\\
\|U(\cdot,t)-U(\cdot,\tau)\|_{L^1(-\infty,\infty)}&\leqslant b|t-\tau|,\qquad \tau,t\in(0,\infty),
\label{align:ch15-1519}\\
\|U(\cdot,t)-\bar U(\cdot,t)\|_{L^1(-\infty,\infty)}&\leqslant a\|U_0(\cdot)-\bar U_0(\cdot)\|_{L^1(-\infty,\infty)},\qquad t\in(0,\infty).
\label{align:ch15-15110}
\end{align}
\end{enumerate}
\end{theorem}

解 $U$ 中的激波满足黏性激波可容性准则，从而满足第八章中所述的一切隐含可容性条件。
当\eqref{align:ch15-1511}的所有特征族要么真非线性、要么线性退化时，$U$ 与由第十三章随机选取法或第十四章波前追踪算法构造的\eqref{align:ch15-1511}、\eqref{align:ch15-1512}解重合。

上述命题的证明综合了多种思想与技巧，将占用本章余下全部篇幅。为便于把握方向，\cref{section:ch15-证明路线图}将提供一幅路线图。

应当指出，估计\eqref{align:ch15-1515}、\eqref{align:ch15-1516}与\eqref{align:ch15-1517}的推导在本质上并不依赖\eqref{align:ch15-1513}具守恒形式这一假设，它同样适用于更一般的系统
\begin{align}
\partial_t U(x,t)+A(U(x,t))\partial_xU(x,t)=\mu\,\partial_x^2U(x,t),\quad x\in(-\infty,\infty),\ t\in(0,\infty),
\label{align:ch15-15111}
\end{align}
只要 $A(U)$ 具有实的互异特征值即可。\eqref{align:ch15-15111}、\eqref{align:ch15-1512}的解族 $\{U_\mu\}$ 的 $\mu\downarrow0$ 极限 $U$ 可解释为
\begin{align}
\partial_tU+A(U)\partial_xU=0
\label{align:ch15-15112}
\end{align}
的一个“弱”解，尽管它未必在分布意义下满足该系统。

\end{document}
